\subsection{Executable extractor contract, cost, and handoff to KBFold}\label{sec:global:handoff}
The proof above replaces the abstract
$\varepsilon_{\rm couple}$ of Chapter~11
by an explicit bounded-replay operation
on the prescribed classical OR blocks.
The following pseudocode fixes which
messages are visible to the \emph{main}
verifier and which are private to the
reduction.

\begin{definition}[Checkpoint extractor pseudocode]\label{def:global:pseudocode}
For each mandatory public OR statement
$Q=(C,S)$, the reduction executes:
\begin{verbatim}
saved = snapshot(prover_before_Q)
(main_bits, main_messages, main_end) =
    run_one_block_with_fresh_uniform_bits(Q)

if not verify_all(main_messages):
    reject_main_and_stop

restore(saved)
repeat at most R times:
    (side_bits, side_messages, _) =
        run_one_block_with_fresh_uniform_bits(Q)
    if verify_all(side_messages) and side_bits != main_bits:
        opening = extract_at_earliest_difference(
                    main_messages, side_messages)
        record_or_compare_with_canonical_registry(opening)
        break
restore(main_end)
continue_main_verifier()
\end{verbatim}
The snapshot includes the private
random tape and complete prover
state. The side executions neither
change the group commitments in the
statement nor publish messages to
the main verifier. The extractor
invokes no irreversible external
or authenticated oracle during
a replayed block.
\end{definition}

\begin{lemma}[Every successful replay extracts from the original statement]\label{lem:global:tag-same}
If the extractor reaches the
\texttt{extract\_at\_earliest\_difference}
line, the extracted opening is for
the original main-path public
statement $Q=(C,S)$.
\end{lemma}
\begin{proof}
Every side run begins at the same
pre-block checkpoint and receives
the same fixed public statement $Q$.
The challenge paths differ only in
fresh inner verifier bits.
By causality, their first messages
are identical up to the earliest
different challenge bit, and the
verifier accepts both mini-proofs
at that index. The OR special
soundness theorem
(Theorem~\ref{thm:realize:or-extract})
returns an opening of the
unchanged group element $C$
to one of the unchanged
candidates in $S$.
\end{proof}

\begin{lemma}[All extracted bit values are registered before use]\label{lem:global:before-use}
On an accepted main run with
successful extraction of every
bitness block, no first-challenge
field variable is defined by an
extraction occurring after $\theta$,
and no second-challenge residual
variable is defined by an extraction
occurring after $z$.
\end{lemma}
\begin{proof}
By Definition~\ref{def:global:epochs},
every mandatory bitness block
for the first-challenge variables
completes in Epoch 0, and all
auxiliary second-challenge bitness
blocks complete in Epoch 1.
The pseudocode always performs
its side replays immediately
after the main block and
\emph{before} continuing to the
next stage. Thus the registered
field values are extracted before
the respective outer sampling
boundaries, as also quantified
in Theorem~\ref{thm:global:chronological}.
\end{proof}

\begin{proposition}[Polynomial overhead for non-negligible cheating advantage]\label{prop:global:polynomial}
Let the cheating advantage be at
least $\varepsilon(\lambda)\ge1/p(\lambda)$
for some polynomial $p$. Choosing
$R=\lceil4Jp\rceil$ makes the
reduction execute at most
\[
 \boxed{O(J^2\kappa p)}
\]
one-bit mini-proof instances,
apart from the cost of restoring
and executing the prover state.
\end{proposition}
\begin{proof}
Lemma~\ref{lem:global:runtime}
bounds the count by
$\kappa J(R+1)$.
With $R=\lceil4Jp\rceil$,
this is at most
$\kappa J(4Jp+2)$,
which has the displayed
polynomial order.
\end{proof}

\begin{definition}[Chapter 13 external interface]\label{def:global:external}
The next chapter receives the
following \emph{fixed contract}:
a binary-field affine verifier
whose public messages and
sequential OR proofs satisfy
Theorem~\ref{thm:global:chronological},
with internal algebraic error
$(\ell+n)/2^n$ and an explicit
computational discrete-log
reduction. It must add a
\emph{separately proven}
authenticated polynomial-oracle
interface that:
\begin{enumerate}[label=\textup{(\roman*)}]
\item binds the original committed
witness/codeword to the extracted
canonical bit registry before
the corresponding challenge;
\item keeps masked opening responses
and Merkle/FRI authentication data
compatible with one globally
simulatable verifier view;
\item preserves KBFold's original
odd-characteristic local
fold/degree and proximity checks,
without transferring a binary-field
characteristic-two identity to
an odd-characteristic word fold.
\end{enumerate}
Its total failure is recorded as
$\varepsilon_{\rm oracle}$ until
a concrete theorem proves it.
\end{definition}

\begin{remark}[No fifth workstream]
All remaining work belongs inside
the locked final structure:
Chapter~13 constructs the full
masked KBFold prover/verifier
and external oracle; Chapter~14
proves or precisely bounds its
end-to-end security; Chapter~15
completes Lean and Rust
certification, complexity and
the final manuscript audit.
No additional numbered chapter
is introduced.
\end{remark}
